Para todo número complexo \(z\) diferente de 0 e 1, definimos a seguinte função
\[  f(z) := \sum \frac 1 {  \log^4 z } \]
onde a soma é sobre todos os ramos do logaritmo complexo.
\begin{itemize}

\item[(a)] Prove que existem dois polinômios \(P\) e \(Q\) tais que \(f(z) = \displaystyle \frac {P(z)}{Q(z)} \) para todo \(z\in\mathbb{C}-\{0,1\}\).

\item[(b)] Prove que para todo \(z\in \mathbb{C}-\{0,1\}\) temos

\[ f(z) = \frac { z^3+4z^2+z}{6(z-1)^4}. \]
\end{itemize}
